A noção de defeitos em teorias de campo quântico (QFTs) tem ganho uma atenção e importância extraordinárias na física e, talvez de modo menos esperado, também na matemática. Seja $M$ uma variedade diferenciável $n$-dimensional na qual vive uma dada QFT $\tau$. Heuristicamente, um defeito de dimensão $d$ pode ser descrito como o dado de uma subvariedade $D\subset M$ de codimensão $(n-d)$, juntamente com alguma descontinuidade de $\tau$ ao longo de $D$ que produz outra teoria de campo quântico. Exemplos paradigmáticos destes fenómenos são os defeitos de linha de Wilson e de ’t Hooft na teoria de Yang--Mills. Dar uma definição geral e uma classificação conceptual destes defeitos é hoje tema de um vasto corpo de trabalho em desenvolvimento.

QFTs que não dependem de quaisquer parâmetros geométricos são chamadas \emph{QFTs topológicas} (TQFTs). Em particular, defeitos cuja interação com a teoria de bulk é independente da classe de isotopia do seu mergulho são chamados defeitos topológicos. Pode então organizar-se uma TQFT topológica com defeitos topológicos num único pacote, chamado teoria topológica de campo quântico com defeitos, ou teoria topológica de campo quântico estratificada. Em linguagem categórica, estas teorias são modeladas por funtores simetricamente monoidais da categoria de cobordismos estratificados com dados de defeito para uma categoria-alvo linear, como a categoria dos espaços vetoriais complexos de dimensão finita. Se enriquecermos estas categorias “para cima”, acrescentando difeomorfismos estratificados, isotopias estratificadas e isotopias de ordem superior, obtemos versões derivadas destas TQFTs; se, por outro lado, as estendermos “para baixo”, permitindo decomposições em peças estratificadas com cantos de profundidade arbitrária, obtemos versões mais computáveis dessas teorias. Esta tese ocupa-se da construção de um tal objeto, numa primeira etapa ao nível de uma $(\oo,1)$-categoria e, de forma prospetiva, ao nível totalmente estendido.

No seu esboço de prova da \emph{Hipótese do Cobordismo}, Lurie introduziu uma primeira encarnação da $(\oo,d)$-categoria completamente estendida de cobordismos framizados como um espaço de Segal completo $d$-vezes iterado. Na secção final desse trabalho, sugere ainda uma definição de uma $(\oo,d)$-categoria de cobordismos com singularidades, baseada numa noção de variedade com singularidades suficientemente geral para capturar defeitos. Em princípio, esse é precisamente o tipo de objeto que aqui pretendemos construir. Contudo, na prática, essa classe de espaços não vinha acompanhada dos fundamentos diferencial-topológicos necessários para uma teoria robusta. A presente tese revê esse programa em três direções complementares: usa um modelo alternativo para categorias superiores, recorre à teoria de espaços estratificados conicamente suaves e às suas estruturas tangenciais, e restringe-se a cobordismos estratificados cujo espaço subjacente é uma variedade topológica.

A tese começa por introduzir, no Capítulo 2, um novo modelo para $(\oo,1)$-categorias, a que chamamos \emph{espaços de Segal intervalares}. Estes são copré-feixes simpliciais na categoria $\Int$, uma subcategoria do poset dos abertos de $[0,1]$, concebida de modo a tornar mais natural a codificação geométrica de composição. Mostramos que a categoria correspondente admite uma estrutura de modelo equivalente de Quillen à estrutura clássica dos espaços de Segal de Rezk. A motivação desta escolha é geométrica: muitos modelos de $(\oo,1)$-categorias são tecnicamente poderosos, mas pouco adequados à construção direta de espaços de módulos geométricos. Uma exceção notável é dada pelas \emph{striation sheaves} de Ayala, Francis e Rozenblyum, desenhadas precisamente para construir $(\oo,1)$-categorias a partir de dados ao nível de conjuntos. Os espaços de Segal intervalares introduzidos nesta tese podem ser vistos como uma alternativa mais próxima da tecnologia tradicional dos espaços de Segal, mas ainda assim suficientemente adaptada a construções geométricas e, por isso, promissora para generalizações a dimensões superiores.

No Capítulo 3 aplicamos este formalismo para construir uma nova versão da $(\oo,1)$-categoria de cobordismos $d$-dimensionais suaves. A abordagem inspira-se em três antecedentes principais: o trabalho de Galatius, Madsen, Tillmann e Weiss sobre o tipo de homotopia da categoria de cobordismos; o programa de Lurie para a hipótese do cobordismo; e a reformulação proposta por Scheimbauer e Calaque. O ponto de partida é que as construções existentes apresentam pequenas falhas técnicas quando reinterpretadas no contexto necessário para uma prova cuidada das condições de Segal e de fibrância. Em particular, a construção de Lurie não satisfaz literalmente uma condição de constância essencial, e a correção posterior de Scheimbauer e Calaque apoia-se numa afirmação problemática sobre a fibrância de certos níveis simpliciais.

Para ultrapassar estas dificuldades, recorremos à teoria de $\oo$-feixes suaves no sítio das variedades diferenciáveis e aos respetivos espaços classificantes, tal como desenvolvidos por Berwick-Evans, Boavida de Brito e Pavlov. A ideia central é construir, a partir de dados geométricos honestos, espaços classificantes de variedades suaves, cobordismos suaves e famílias suaves de composições de cobordismos. Mais concretamente, definimos um bifuntor $\widetilde{\bord}_{d,V}(-,-)$ em $\man^{\op}\times \Int$ que, para cada intervalo $T$, produz um feixe de conjuntos sobre variedades suaves e codifica famílias de variedades $d$-dimensionais mergulhadas num espaço euclidiano auxiliar $V$. Tomando a sua resolução por concordâncias e restringindo depois ao longo da variedade co-simplicial padrão, obtemos um conjunto simplicial $\bord_{d,V}$ que funciona como espaço classificante de variedades $(d-1)$-dimensionais, de cobordismos $d$-dimensionais e das suas composições. A parte técnica central do capítulo consiste em mostrar que este pré-espaço intervalar satisfaz as condições que o tornam num espaço de Segal intervalar e que a sua categoria de homotopia 1-truncada recupera a categoria de cobordismos esperada.

No Capítulo 4 passamos à geometria estratificada propriamente dita. Introduzimos a teoria dos \emph{espaços estratificados conicamente suaves}, desenvolvida por Ayala, Francis, Tanaka e Rozenblyum, como o contexto geométrico adequado para modelar defeitos. Esta escolha é motivada por duas exigências fundamentais: trabalhar com famílias de estratificações e formular noções suficientemente flexíveis de estrutura tangencial em contexto singular. O capítulo desenvolve o ferramental diferencial-topológico necessário para os capítulos seguintes: a $(\oo,1)$-categoria de espaços estratificados conicamente suaves, a categoria de caminhos de saída e o fibrado tangente construtível. Para além disso, demonstramos resultados novos sobre a contratilidade de certos espaços de mergulhos de espaços estratificados em espaços euclidianos, bem como um análogo estratificado do teorema de Ehresmann. Estes resultados são essenciais para transportar para o mundo estratificado parte da infraestrutura já conhecida no caso liso.

O Capítulo 5 contém o núcleo geométrico da tese: a construção de uma $(\oo,1)$-categoria de cobordismos $d$-dimensionais estratificados conicamente suaves. Começamos por introduzir a noção de variedade topológica estratificada conicamente suave, modelada como uma $\cB$-variedade estratificada para uma adequada $(\oo,1)$-categoria de básicos $\Bsc_{d}^{\cS_{d-1}}$, cujas ligações são esferas estratificadas das codimensões relevantes. Esta restrição a espaços subjacentes topológicos permite manter um equilíbrio entre generalidade geométrica e tratabilidade técnica. Construímos então um pré-espaço intervalar de cobordismos estratificados, novamente através da resolução por concordâncias de $\oo$-feixes suaves e da restrição ao longo da variedade co-simplicial padrão, e provamos que ele satisfaz a condição de Segal intervalar.

Em seguida, dotamos este espaço de uma estrutura simetricamente monoidal compatível com a união disjunta, obtendo assim o candidato natural à categoria de cobordismos estratificados necessária para formular TQFTs estratificadas derivadas. Um aspeto importante desta construção é a incorporação de estruturas tangenciais no contexto estratificado. Em particular, introduzimos a noção de \emph{dado de defeito discreto orientado} através de fibrados à direita discretos sobre a categoria de básicos escolhida, com o objetivo de aproximar o formalismo geométrico desenvolvido nesta tese dos dados de defeito considerados por Carqueville e colaboradores. Como primeiro teste desta perspetiva, provamos um resultado de comparação no caso orientado estratificado de dimensão $2$. O capítulo culmina com a definição de uma teoria topológica de campo quântico estratificada derivada como um funtor simétrico monoidal da $(\oo,1)$-categoria simétrica monoidal de cobordismos estratificados construída na tese para uma $(\oo,1)$-categoria simétrica monoidal alvo $\cC$.

O Capítulo 6 tem natureza prospetiva e delineia duas direções de desenvolvimento futuro. Na primeira, seguindo ideias não publicadas de Manuel Araújo e Pedro Boavida de Brito, introduzimos uma generalização de dimensão superior da categoria $\Int$, denotada $\cube_d$, com vista à formulação de um análogo intervalar para $(\oo,d)$-categorias. Argumentamos que copré-feixes simpliciais em $\cube_d$ satisfazendo análogos superiores das condições de Segal e de completude deverão apresentar a $(\oo,1)$-categoria $\Cat_{(\oo,d)}$, e esboçamos como este formalismo poderá ser usado para construir um $\cube_d$-espaço de cobordismos estratificados. Na segunda direção regressamos ao contexto não estendido e indicamos como álgebras de fatorização construtíveis e homologia de fatorização estratificada podem servir de mecanismo de produção de TQFTs estratificadas com defeitos, generalizando para o contexto estratificado e derivado ideias presentes na abordagem de Calaque e Scheimbauer.

Em síntese, a tese propõe um novo quadro para a construção de categorias superiores geométricas adaptadas ao estudo de cobordismos estratificados e de teorias topológicas de campo quântico com defeitos. A introdução dos espaços de Segal intervalares, a nova construção da categoria de cobordismos lisos via $\oo$-feixes suaves, e a passagem para o contexto dos espaços estratificados conicamente suaves constituem os três pilares principais deste programa. O resultado é uma construção rigorosa de uma $(\oo,1)$-categoria de cobordismos estratificados suficientemente rica para servir de domínio a TQFTs estratificadas derivadas, bem como um esboço das extensões futuras necessárias para tratar a versão totalmente estendida do problema.

\bigskip

\begin{samepage}
\noindent\textbf{Palavras-chave:}
\begin{itemize}[nosep]
    \item espaços estratificados conicamente suaves
    \item bordismos estratificados
    \item espaços de Segal intervalares
    \item categorias superiores
    \item teorias de campo quântico topológicas
\end{itemize}
\end{samepage}
